\section*{第 \ref{ch:g7:states-of-matter} 章 --- 物态与物态变化}

\begin{solution}{exo:g7:states-of-matter:1}
固体：分子整齐堆叠，原地振动。
液体：分子相互接触但杂乱无序，彼此滑过；
气体：分子相距很远，向各个方向快速飞动。
\end{solution}

\begin{solution}{exo:g7:states-of-matter:2}
在固体中，每个分子都守住自己在堆叠中的位置，
因此整体保持形状。在液体中，分子仍然相互接触，
却可以自由地彼此滑过：容器中的分子群下滑、填满容器，
直到再也无法下滑 --- 形成水平的液面。
\end{solution}

\begin{solution}{exo:g7:states-of-matter:3}
空气主要是飞行的分子之间的空隙：推压会把
这些飞行者挤得更近（它们的撞击会反抗 --- 好似弹簧）。
水的分子已经相互接触，没有空隙可以挤掉，
因此活塞会碰到一堵墙。
\end{solution}

\begin{solution}{exo:g7:states-of-matter:4}
气体分子相距很远，之间是大片空旷 ---
因此整间房的气体都能被压缩到一起。水的分子
已经相互接触，是不可压缩的：任何泵都无法将它压得更紧。
\end{solution}

\begin{solution}{exo:g7:states-of-matter:5}
加热使堆叠中的分子振动得越来越剧烈；
到达熔点时，振动压倒了堆叠的束缚，
行列瓦解，分子开始滑动：
固体便变成了液体。
\end{solution}

\begin{solution}{exo:g7:states-of-matter:6}
物态变化只是重新排列同样的分子 --- 分子既不会产生，
也不会消失，因此质量不会改变。要使质量减少，
分子就必须消失（或逃出容器 ---
这是未盖瓶盖的瓶子造成的假象）。
\end{solution}

\begin{solution}{exo:g7:states-of-matter:7}
分子运动的活跃程度。热的物体与冷的物体接触时，
活跃的分子撞击迟缓的分子，每次碰撞都传递一点活力，
活跃者平静下来，迟缓者随之加快 ---
直至两群分子步调一致，温度相等。
\end{solution}

\begin{solution}{exo:g7:states-of-matter:8}
香水分子飞出，与空气分子相撞、反弹、游荡，
一步步曲折地扩散到整个房间。微风会整团地裹挟它们：
顺风而行，胜过在人群中
跌跌撞撞地穿行。
\end{solution}

\begin{solution}{exo:g7:states-of-matter:9}
只有\emph{最快}的分子才能挣脱液体的束缚，
因此蒸发是对活跃程度的一种“征税”：逃逸者带走的
超过了它们应有的份额，留下的分子群平均而言
变慢了 --- 也就是变凉了。汗水和湿手指
正是靠带走最活跃的分子而降温。
\end{solution}

\begin{solution}{exo:g7:states-of-matter:10}
阳光下，被困的分子飞得更快，撞击皮肤更猛、更频繁：
如冰雹般的撞击使气球膨胀，直到绷紧的表皮与之平衡。
寒冷中，同样的分子变慢，撞击减弱，
外界空气的压力便把
气球压小。
\end{solution}

\begin{solution}{exo:g7:states-of-matter:11}
水结冰时，分子锁定成一种异常疏松、开放的堆叠，
空隙比流动的分子群中更多，因此同样多的分子
占据更大的体积：密度小于液体，所以冰会浮起。
若水是普通液体，冰会紧密堆叠，形成时就会下沉，
池塘会从底部向上结冰 --- 那就不会有避风的水
留给鱼了。
\end{solution}

\begin{solution}{exo:g7:states-of-matter:12}
没有任何东西搅动水，墨水却在移动：液体自身的分子
处于永不停歇的运动之中，推挤着墨水分子，
一步步随机地穿过分子群 --- 水和墨水两种液体
每时每刻都充满运动，即使玻璃杯看起来
完全静止。温度升高，每次推挤都更猛，扩散也就更快。
耐心的墨水由此使看不见的舞动变得可见 ---
这一证据使那些怀疑分子存在的人信服。
\end{solution}

\begin{solution}{pb:g7:states-of-matter:1}
\textbf{1.} 整齐排成若干行，肩并肩，每人原地
抖动 --- 谁都不许换位置或离开自己的位置。
\textbf{2.} 队列瓦解：学生仍然肩并肩
（相互接触！）但现在可以来回穿插、彼此滑过。他们仍须
与人群保持接触 --- 谁也不许自由跑开。
\textbf{3.} 学生四散奔跑，快速地沿直线前进，直到
相撞或撞到墙壁，然后不停地反弹 --- 既然没有东西
把他们连在一起，只有墙壁能让他们停下，人群
占满了体育馆的每个角落。
\textbf{4.} 是的 --- 每一幕都是同样的三十人：质量
在物态变化中保持不变，这是舞台上演示的。
\textbf{5.} 气体一幕允许这样做 --- 跑步者挤进
体育馆的一半（并且把牧羊人撞得更猛）。液体那一团不能缩小：
学生已经相互接触。演示：气体可被压缩，
液体则不能。
\textbf{6.} 在各自位置上越抖越剧烈；
抖动冲破队列的那一刻，这一幕便成了液体，
学生开始彼此滑过。
\textbf{7.} 刻度盘的能量用于打破堆叠的束缚 ---
一排一排地打破 --- 而不用于加快运动：温度计评判员
必须报告，在最后一排瓦解之前，平均活跃程度\emph{保持不变}。
平台期，就此演示出来。
\textbf{8.} 只有\emph{最快}的学生 --- 只有他们
能够挣脱。其余人群的平均活跃程度下降：液体
通过蒸发使自己变凉，
与汗水和湿手指的凉意如出一辙。
\textbf{9.} 演员不断与绳子相撞，并将它向外推；
速度更快的演员推得更猛、更频繁，圆环便膨胀起来：
正如阳光下的气球。
\textbf{10.} “以\emph{疏松}的次序堆叠 --- 手臂的距离，
空隙多”：水的固体比液体占据更多空间，因此
冰一幕中的木筏会浮在液体一幕的人群之上。
\textbf{11.} 蓝色衬衫被推挤 --- 一步步随机地
撞过滑动的人群，直到游荡遍各处：墨水在扩散。
把刻度盘调高，每次推挤都更快，
游荡也就更快。
\textbf{12.} 例如：“今晚的演出只呈现一个观点：一切物质
都是运动中的分子。诚然，两条旧定律也随之而来 ---
质量在每一次物态变化中保持不变，以及熔化时的
平台期。我们无法演出的是尺度：真实的演员
比任何能被看见的东西都小，
数量比世界上所有的观众还多。”
\end{solution}
